\documentclass[11pt,a4paper]{article}

\usepackage[utf8]{inputenc}
\usepackage[T1]{fontenc}
\usepackage[spanish,es-noshorthands]{babel}
\usepackage{xcolor}
\usepackage[most]{tcolorbox}
\usepackage{listings}
\usepackage{etoolbox}
\usepackage{fancyhdr}
\usepackage{lastpage}
\usepackage{enumitem}
\usepackage[hidelinks]{hyperref}

\setlength{\headheight}{30pt}
\pagestyle{fancy}
\fancyhf{}
\renewcommand{\headrulewidth}{0pt}
\fancyhead[R]{%
  \fontsize{8}{9.6}\selectfont
  Licenciatura en Ciencias de la Computación\\
  Sistemas Operativos\\[2pt]
  \rule{4.2cm}{0.5pt}
}
\fancyfoot[C]{\thepage\ / \pageref{LastPage}}

\newtcolorbox{trackbox}[1]{
  colback=green!7!white, colframe=green!45!black,
  arc=2mm, boxrule=0.6pt, left=8pt, right=8pt, top=6pt, bottom=6pt,
  title=TRACK -- #1, fonttitle=\bfseries, coltitle=green!35!black,
  colbacktitle=green!7!white, boxsep=1pt, breakable
}

\newtcolorbox{extrabox}{
  colback=orange!8!white, colframe=orange!55!white,
  arc=2mm, boxrule=0.6pt, left=8pt, right=8pt, top=6pt, bottom=6pt,
  title=EXTRA, fonttitle=\bfseries, coltitle=orange!55!black,
  colbacktitle=orange!8!white, boxsep=1pt, breakable
}

\newtcolorbox{bonusbox}{
  colback=red!7!white, colframe=red!55!black,
  arc=2mm, boxrule=0.6pt, left=8pt, right=8pt, top=6pt, bottom=6pt,
  title=BONUS -- POR PUNTOS, fonttitle=\bfseries, coltitle=red!55!black,
  colbacktitle=red!7!white, boxsep=1pt, breakable
}

\lstdefinestyle{codestyle}{
  backgroundcolor=\color{black!5}, rulecolor=\color{black!25},
  frame=single, framerule=0.5pt, basicstyle=\ttfamily\small,
  breaklines=true, showstringspaces=false, tabsize=4,
  keywordstyle=\color{blue!60!black}\bfseries,
  commentstyle=\color{green!40!black},
  stringstyle=\color{orange!70!black},
  columns=fullflexible
}
\lstset{style=codestyle}
\BeforeBeginEnvironment{lstlisting}{%
  \par\addvspace{0.2\baselineskip}\noindent
  \begin{minipage}{\linewidth}}
\AfterEndEnvironment{lstlisting}{%
  \end{minipage}\par\addvspace{0.2\baselineskip}}
\newcommand{\inputcode}[2][]{%
  \par\addvspace{0.2\baselineskip}\noindent
  \begin{minipage}{\linewidth}
    \lstinputlisting[#1]{#2}
  \end{minipage}\par\addvspace{0.2\baselineskip}}
\setlist{nosep}

\title{Trabajo Práctico 3 -- Ejercicios extra}
\author{}
\date{}

\begin{document}
\renewcommand{\contentsname}{Tracks}
\maketitle
\thispagestyle{fancy}
\begin{center}
SISTEMAS OPERATIVOS
\end{center}
\newpage

\tableofcontents
\newpage

Estos tracks amplían el práctico de memoria con experimentos comparables.
Registrá arquitectura, versión de \texttt{gcc} y opciones de compilación; las
direcciones concretas no son portables, así que explicá patrones, permisos,
tamaños y relaciones.

Los tracks son independientes entre sí: elegí uno después de completar el
práctico regular y no necesitás resolver otro track para empezar. Las cajas
\textbf{EXTRA} proponen una extensión opcional del track; las cajas
\textbf{BONUS -- POR PUNTOS} plantean una implementación más profunda. El
práctico regular no requiere un informe; para obtener puntos de bonus, entregá
una explicación breve con capturas y quedate disponible para una demostración.

\newpage
\section[Páginas bajo demanda y memoria residente (7/10)]{Páginas bajo demanda y memoria residente\\[-2pt](7/10)}

\begin{trackbox}{7/10}
\textbf{Descripción:} diferenciar el espacio virtual reservado de las páginas
tocadas y residentes.
\end{trackbox}

Escribí un programa \texttt{demanda.c} que demuestre la paginación bajo
demanda (\emph{demand paging}). Observá los campos \texttt{Vm*} de
\texttt{/proc/<pid>/status}; provocá accesos a páginas nuevas y registrá los
fallos \texttt{minflt}/\texttt{majflt} de \texttt{/proc/<pid>/stat}. Relacioná
esos fallos con los cambios que observes en
\texttt{/proc/<pid>/smaps\_rollup} y explicá el resultado.

\begin{lstlisting}[language=bash]
$ gcc -Wall -Wextra -O0 -g demanda.c -o demanda
$ ./demanda &
$ pid=$!
$ grep '^Vm' /proc/$pid/status
$ awk '{printf "minflt=%s majflt=%s\n", $10, $12}' /proc/$pid/stat
$ cat /proc/$pid/smaps_rollup
$ wait $pid
\end{lstlisting}

\begin{extrabox}
Formulá una hipótesis antes de tocar las páginas. Explicá por qué
\texttt{malloc()} puede aumentar \texttt{VmSize} sin aumentar inmediatamente
\texttt{VmRSS}, y qué evento provoca la primera escritura en cada página.
\end{extrabox}

\newpage
\section{\texttt{mmap()}, permisos y regiones anónimas (7/10)}

\begin{trackbox}{7/10}
\textbf{Descripción:} crear y modificar mapeos explícitos para relacionar
llamadas al sistema, permisos y entradas de \texttt{maps}.
\end{trackbox}

\begin{enumerate}[leftmargin=*]
  \item Reservar una región anónima con \texttt{mmap()} usando
  \texttt{PROT\_READ} $|$ \texttt{PROT\_WRITE} y
  \texttt{MAP\_PRIVATE} $|$ \texttt{MAP\_ANONYMOUS}.
  \item Imprimir la dirección, escribir en la primera y última página y
  localizar la región en \texttt{/proc/<pid>/maps}.
  \item Usar \texttt{mprotect()} para quitar escritura y volver a inspeccionar.
  \item Mapear un archivo real con \texttt{MAP\_PRIVATE}, modificar un byte
  mediante el mapeo y comprobar que el archivo no cambia. Explicá qué significa
  que la modificación sea privada y contrastá, opcionalmente, con
  \texttt{MAP\_SHARED}.
  \item Liberar las regiones con \texttt{munmap()} y comprobar que desaparecen.
\end{enumerate}

\begin{lstlisting}[language=C]
void *region = mmap(NULL, 4 * 4096,
                    PROT_READ | PROT_WRITE,
                    MAP_PRIVATE | MAP_ANONYMOUS, -1, 0);
if (region == MAP_FAILED)
    perror("mmap");

if (mprotect(region, 4 * 4096, PROT_READ) == -1)
    perror("mprotect");
\end{lstlisting}

Usá \texttt{pmap -x <pid>} como vista alternativa y compará sus columnas con
\texttt{maps}. Investigá por qué una región privada puede compartir páginas
físicamente mediante \emph{copy-on-write}.

\newpage
\section{Intercambio y presión de memoria (2/10)}

\begin{trackbox}{2/10}
\textbf{Descripción:} activar temporalmente un archivo de intercambio
(\emph{swapfile}) y observar cómo aparece en las herramientas del sistema.
\end{trackbox}

Trabajá en una máquina de prueba y usá un archivo descartable. Necesitás
privilegios para ejecutar \texttt{swapon}; no modifiques \texttt{/etc/fstab},
no uses \texttt{swapoff -a} y no desactives particiones o archivos de
intercambio que ya estuvieran activos. El archivo debe estar en un sistema de
archivos que admita swapfiles.

Desde un directorio con espacio libre, creá y activá un archivo de 256 MiB:

\begin{lstlisting}[language=bash]
$ swapon --show
$ df -h .
$ fallocate -l 256M swapfile-demo
$ chmod 600 swapfile-demo
$ sudo mkswap swapfile-demo
$ sudo swapon swapfile-demo
$ swapon --show
$ cat /proc/swaps
$ free -h
\end{lstlisting}

Si \texttt{fallocate} no está disponible o el sistema de archivos rechaza el
archivo, investigá una alternativa que cree un archivo completamente
asignado, por ejemplo:

\begin{lstlisting}[language=bash]
$ dd if=/dev/zero of=swapfile-demo bs=1M count=256 status=progress
\end{lstlisting}

Mientras ejecutás una carga de memoria
moderada, observá \texttt{vmstat 1}, \texttt{free -h} y
\texttt{/proc/meminfo}. Distinguí entre tener swap habilitado y que el núcleo
efectivamente mueva páginas hacia él; no es necesario agotar la memoria de la
máquina para realizar la observación.

Al terminar, desactivá solamente el archivo de la práctica, comprobá que ya no
aparece y recién después borrálo:

\begin{lstlisting}[language=bash]
$ sudo swapoff swapfile-demo
$ swapon --show
$ rm swapfile-demo
\end{lstlisting}

Si \texttt{swapoff} falla porque el archivo todavía está en uso, no lo borres.

\newpage
\section{Bibliotecas compartidas y resolución de símbolos (5/10)}

\begin{trackbox}{5/10}
\textbf{Descripción:} investigar cómo el cargador encuentra una biblioteca
compartida y cuándo resuelve un símbolo del ejecutable.
\end{trackbox}

Partí de los archivos \texttt{saludo.c} y \texttt{main.c} del TP3 regular:
\texttt{saludo.c} define la función de la biblioteca y \texttt{main.c} la
invoca. Construí una variante dinámica sin una ruta de búsqueda incrustada,
para que puedas investigar cómo el cargador encuentra \texttt{libsaludo.so}:

\begin{lstlisting}[language=bash]
$ gcc -Wall -Wextra -fPIC -shared saludo.c \
    -Wl,-soname,libsaludo.so -o libsaludo.so
$ gcc -Wall -Wextra main.c -L. -lsaludo -o saludo
$ ldd ./saludo
$ readelf -dW ./saludo
$ readelf -dW ./libsaludo.so
$ readelf -rW ./saludo
$ LD_LIBRARY_PATH=. ./saludo
$ LD_DEBUG=libs,bindings LD_LIBRARY_PATH=. ./saludo 2>&1 | less
\end{lstlisting}

Prepará una segunda copia de \texttt{libsaludo.so} en un directorio diferente,
con un mensaje distinguible, y compará:

\begin{lstlisting}[language=bash]
$ LD_LIBRARY_PATH=alternativa:. ./saludo
$ LD_LIBRARY_PATH=.:alternativa ./saludo
\end{lstlisting}

Explicá qué biblioteca se elige en cada caso; usá \texttt{LD\_DEBUG} para
justificar la búsqueda y \texttt{readelf -dW} y \texttt{readelf -rW} para
relacionarla con \texttt{SONAME}, \texttt{NEEDED} y las reubicaciones. Compará
el enlace perezoso con \texttt{LD\_BIND\_NOW=1} y describí cuándo se resuelve
\texttt{saludar}. No copies direcciones.

\newpage
\section{Inspección de un plugin VST (4/10)}

\begin{trackbox}{4/10}
\textbf{Descripción:} inspeccionar un plugin VST (\emph{Virtual Studio
Technology}) real como una biblioteca compartida, sin ejecutarlo.
\end{trackbox}

Elegí un plugin VST o VST3 para Linux de una fuente confiable. No hace falta
instalar un DAW (\emph{Digital Audio Workstation}) ni producir audio. Si el
plugin tiene formato \texttt{.vst3}, localizá dentro del paquete la biblioteca
correspondiente a Linux; si es un archivo \texttt{.so}, usalo directamente.

Inspeccioná la biblioteca con las herramientas del track anterior:

\begin{enumerate}[leftmargin=*]
  \item Usá \texttt{file} para identificar el formato y la arquitectura.
  \item Usá \texttt{ldd} para listar sus dependencias y comprobar si falta
  alguna.
  \item Usá \texttt{readelf -hW} y \texttt{readelf -dW} para observar el tipo
  de archivo, el intérprete y las bibliotecas requeridas.
  \item Usá \texttt{nm -D --defined-only} para observar algunos símbolos
  exportados.
  \item Explicá qué tendría que hacer un programa anfitrión para usar esta
  biblioteca y qué riesgo existe al cargar código de terceros.
\end{enumerate}

Registrá el nombre y la versión del plugin, la arquitectura, sus dependencias
y tres observaciones de los comandos. No copies direcciones: relacioná los
resultados con la carga de bibliotecas compartidas.

\begin{bonusbox}
Implementá un cargador de plugins usando \texttt{dlopen()}, \texttt{dlsym()}
y \texttt{dlclose()}.

Definí una interfaz binaria mínima (ABI, \emph{Application Binary Interface}),
por ejemplo una función \texttt{int plugin\_run(const char *argumento)}. Creá
dos plugins \texttt{.so} que implementen esa función y un programa
\texttt{runner} que reciba por línea de comandos la ruta del plugin, lo cargue
con \texttt{RTLD\_NOW | RTLD\_LOCAL}, busque el símbolo y lo ejecute.

El programa debe informar errores de \texttt{dlopen()} y \texttt{dlsym()}
mediante \texttt{dlerror()}, y cerrar la biblioteca con \texttt{dlclose()}.

Compará este programa con otro enlazado directamente contra una biblioteca
mediante \texttt{-l}. Usá \texttt{readelf -dW},
\texttt{LD\_DEBUG=libs,bindings} y \texttt{/proc/<pid>/maps} para explicar:
\begin{itemize}[leftmargin=*]
  \item por qué el plugin cargado con \texttt{dlopen()} no aparece como
  dependencia \texttt{NEEDED} inicial;
  \item cuándo se resuelve el símbolo;
  \item qué diferencia hay entre \texttt{RTLD\_LOCAL} y una biblioteca
  enlazada normalmente;
  \item qué puede fallar si el plugin no existe, no exporta el símbolo o usa
  una interfaz incompatible;
  \item qué riesgos tiene cargar bibliotecas proporcionadas por terceros.
\end{itemize}
\end{bonusbox}

\newpage
\section{Diagnóstico de errores de memoria (6/10)}

\begin{trackbox}{6/10}
\textbf{Descripción:} provocar errores controlados y comparar Valgrind,
AddressSanitizer y un depurador.
\end{trackbox}

Usá este programa sólo en un entorno de prueba:

\inputcode[language=C]{../examples/tp3-extra/errores.c}

\begin{enumerate}[leftmargin=*]
  \item Ejecutarlo sin herramientas y observar si el error es determinista.
  \item Usar \texttt{valgrind --leak-check=full ./errores}.
  \item Compilar con \texttt{-fsanitize=address -fno-omit-frame-pointer -g}
  y repetir.
  \item Ejecutarlo bajo \texttt{gdb}: detenerse en \texttt{main}, avanzar con
  \texttt{next} y obtener un \texttt{backtrace} cuando aparezca el error.
  \item Comparar pila, ubicación y momento en que cada herramienta detiene el
  programa.
  \item Observar las bibliotecas del runtime en
  \texttt{/proc/<pid>/maps} mientras el programa está detenido.
\end{enumerate}

\begin{lstlisting}[language=bash]
$ gcc -Wall -Wextra -g errores.c -o errores
$ valgrind --leak-check=full ./errores
$ gcc -Wall -Wextra -g -fsanitize=address \
    -fno-omit-frame-pointer errores.c -o errores-asan
$ ./errores-asan
\end{lstlisting}

\begin{lstlisting}[language=bash]
$ gdb ./errores
(gdb) break main
(gdb) run
(gdb) next
(gdb) continue
(gdb) backtrace
\end{lstlisting}

\begin{bonusbox}
Construí una versión corregida con tres modos de error seleccionables
(desborde, uso después de \texttt{free} y fuga) y un script que ejecute cada
modo con Valgrind y AddressSanitizer. Compará los diagnósticos y explicá por
qué un fallo de segmentación no necesariamente señala la línea donde terminó
el proceso. Incluí capturas y dejá preparada una demostración.
\end{bonusbox}

\end{document}
